\section{Nueva frontera I: IA generativa y modelos de difusión}

\subsection{Trayectoria del desarrollo del modelado generativo}

En la Lecture 4, el curso presentó dos modelos generativos clásicos: el VAE (Variational Autoencoder) y la GAN (Generative Adversarial Network). Aunque estos dos tipos de modelos lograron, de manera pionera, la capacidad de generar nuevas muestras a partir de datos, enfrentan varias limitaciones clave:

\begin{figure}[H]
\centering
\includegraphics[width=0.85\textwidth]{slides-images/slide-48.jpg}
\caption{Limitaciones y retos del VAE/GAN\protect\footnotemark}
\end{figure}
\footnotetext{Diapositiva 48.}

\begin{knowledgebox}{Principales problemas del VAE y la GAN}
\textbf{Limitaciones (Limitations)}:
\begin{itemize}
    \item \textbf{Mode collapse}: el modelo tiende a generar únicamente muestras ``promedio'', perdiendo diversidad
    \item \textbf{Dificultad para generar fuera de la distribución}: es difícil generar muestras novedosas fuera de la distribución de entrenamiento
    \item \textbf{Dificultad de entrenamiento}: en particular, el proceso de entrenamiento adversario de la GAN es sumamente inestable
\end{itemize}
\textbf{Retos (Challenges)}: estabilidad, eficiencia, calidad de generación, novedad
\end{knowledgebox}

\subsection{Idea central de los Diffusion Models}

Los Diffusion Models proponen un paradigma de generación completamente distinto al del VAE/GAN:

\begin{figure}[H]
\centering
\includegraphics[width=0.85\textwidth]{slides-images/slide-50.jpg}
\caption{Generación ``one-shot'' del VAE/GAN vs. generación por eliminación iterativa de ruido de Diffusion\protect\footnotemark}
\end{figure}
\footnotetext{Diapositiva 50.}

\begin{importantbox}{Innovación clave de los Diffusion Models}
\begin{itemize}
    \item \textbf{VAE/GAN}: genera la muestra completa \textbf{de una sola vez} (one-shot) a partir de una variable latente de baja dimensión
    \item \textbf{Diffusion Models}: genera la muestra eliminando el ruido \textbf{de manera iterativa} (iteratively) y por pasos
\end{itemize}
Esta estrategia iterativa descompone una tarea de generación difícil en una serie de pasos simples de eliminación de ruido, lo que mejora notablemente la calidad de generación.
\end{importantbox}

\subsection{Proceso de adición de ruido hacia adelante (Forward Noising)}

El primer paso de un modelo de Diffusion es construir los datos de entrenamiento mediante un proceso de adición de ruido hacia adelante (Forward Noising Process).

\begin{figure}[H]
\centering
\includegraphics[width=0.85\textwidth]{slides-images/slide-51.jpg}
\caption{Proceso de difusión: adición de ruido hacia adelante y eliminación de ruido hacia atrás\protect\footnotemark}
\end{figure}
\footnotetext{Diapositiva 51. Cita a Sohl-Dickstein y otros, ICML 2015; Ho y otros, NeurIPS 2020.}

\begin{figure}[H]
\centering
\includegraphics[width=0.85\textwidth]{slides-images/slide-53.jpg}
\caption{Proceso de adición de ruido hacia adelante: adición progresiva de ruido\protect\footnotemark}
\end{figure}
\footnotetext{Diapositiva 53.}

Pasos clave del proceso hacia adelante:
\begin{enumerate}
    \item Se parte de una imagen de entrenamiento $x_0$
    \item Se toma una muestra de un patrón de ruido aleatorio
    \item Se aumenta gradualmente la proporción de ruido según un programa (schedule) predefinido:
    \begin{itemize}
        \item $T=0$: 100\% imagen, 0\% ruido
        \item $T=1$: 75\% imagen, 25\% ruido
        \item $T=2$: 50\% imagen, 50\% ruido
        \item $T=3$: 25\% imagen, 75\% ruido
        \item $T=4$: 0\% imagen, 100\% ruido (ruido puramente aleatorio)
    \end{itemize}
\end{enumerate}

\begin{knowledgebox}{El proceso hacia adelante no requiere entrenamiento}
El proceso de adición de ruido hacia adelante es completamente determinista (dado un schedule de ruido) y no involucra ningún parámetro entrenable. Su único propósito es generar pares de entrenamiento para el proceso de eliminación de ruido hacia atrás: cada par $(x_t, x_{t-1})$ constituye una muestra de entrenamiento.
\end{knowledgebox}

\subsection{Proceso de eliminación de ruido hacia atrás (Reverse Denoising)}

La tarea central del entrenamiento es que la red neuronal aprenda el proceso inverso: dada la imagen con ruido en el instante $T$, estimar la imagen más limpia del instante $T-1$.

\begin{figure}[H]
\centering
\includegraphics[width=0.85\textwidth]{slides-images/slide-54.jpg}
\caption{Eliminación de ruido hacia atrás: aprender a estimar $T-1$ a partir de $T$\protect\footnotemark}
\end{figure}
\footnotetext{Diapositiva 54. Cita a Ho y otros, NeurIPS 2020.}

\begin{importantbox}{Objetivo de entrenamiento de Diffusion}
Dada la imagen con ruido $x_T$ en el instante $T$, se entrena una red neuronal $f_\theta$ para predecir la imagen $x_{T-1}$ del instante $T-1$:
$$
\hat{x}_{T-1} = f_\theta(x_T, T)
$$
La pérdida de entrenamiento es la diferencia entre la imagen predicha y la imagen real (habitualmente el MSE, o el MSE del ruido predicho). Esta red comparte parámetros en todos los pasos de tiempo, y el paso de tiempo $T$ se usa como entrada de condicionamiento.
\end{importantbox}

\subsection{Muestreo para generar imágenes nuevas}

Una vez terminado el entrenamiento, el proceso de generar una imagen nueva consiste en partir de ruido puro y aplicar repetidamente la red de eliminación de ruido:

\begin{figure}[H]
\centering
\includegraphics[width=0.85\textwidth]{slides-images/slide-61.jpg}
\caption{Proceso de muestreo: de ruido puro a una imagen nítida, paso a paso\protect\footnotemark}
\end{figure}
\footnotetext{Diapositiva 61.}

Proceso de muestreo: $x_T \xrightarrow{f_\theta} x_{T-1} \xrightarrow{f_\theta} x_{T-2} \xrightarrow{f_\theta} \cdots \xrightarrow{f_\theta} x_0$

Cada paso de eliminación de ruido solo necesita retirar ``una pequeña parte'' del ruido, lo cual es mucho más sencillo que generar la imagen completa de una sola vez; esta es precisamente la razón fundamental de la calidad superior de generación de los Diffusion Models.

\subsection{De la imagen al lenguaje natural: Text-to-Image}

Una aplicación importante de los Diffusion Models es su combinación con el procesamiento de lenguaje natural: la generación de imágenes a partir de texto (Text-to-Image Generation).

\begin{figure}[H]
\centering
\includegraphics[width=0.85\textwidth]{slides-images/slide-64.jpg}
\caption{Generación de imágenes a partir de descripciones en lenguaje natural\protect\footnotemark}
\end{figure}
\footnotetext{Diapositiva 64. Cita a Ramesh y otros, arXiv 2022.}

\begin{figure}[H]
\centering
\includegraphics[width=0.85\textwidth]{slides-images/slide-65.jpg}
\caption{Ejemplos diversos de generación Text-to-Image\protect\footnotemark}
\end{figure}
\footnotetext{Diapositiva 65. Cita a OpenAI, Ramesh y otros, arXiv 2022.}

Los sistemas Text-to-Image dominantes actuales (como DALL-E, Stable Diffusion, Midjourney) usan mayoritariamente un Diffusion Model como red troncal de generación de imágenes. El texto introducido por el usuario se transforma, mediante un codificador de lenguaje, en un vector de condicionamiento que guía el proceso de eliminación de ruido para que genere la imagen en la dirección que corresponde a la descripción textual.

\subsection{Más allá de la imagen: diseño molecular}

\begin{figure}[H]
\centering
\includegraphics[width=0.85\textwidth]{slides-images/slide-66.jpg}
\caption{Aplicación de los Diffusion Models en el diseño molecular y la generación de proteínas\protect\footnotemark}
\end{figure}
\footnotetext{Diapositiva 66.}

Las aplicaciones de los Diffusion Models van mucho más allá de la generación de imágenes. En los campos de la química y la biología, se emplea el mismo paradigma de ``del ruido a la estructura'' para:
\begin{itemize}
    \item \textbf{Generación de moléculas}: generar de manera progresiva estructuras de moléculas de fármacos a partir de ruido 3D
    \item \textbf{Diseño de proteínas}: generar nuevas estructuras de proteínas con funciones específicas
\end{itemize}

\begin{knowledgebox}{Generalidad de Diffusion}
El marco de ``adición de ruido $\rightarrow$ eliminación de ruido'' de los Diffusion Models tiene una generalidad muy alta: basta con que los datos puedan ser ``corrompidos'' y ``restaurados'' de manera progresiva para que este marco pueda aplicarse. Esto lo convierte en uno de los paradigmas de modelado generativo más activos en la actualidad, con aplicaciones que abarcan imágenes, video, audio, estructuras 3D, moléculas y muchos otros campos.
\end{knowledgebox}

\subsection{Resumen de la sección}

Los Diffusion Models, al descomponer la tarea de generación en múltiples pasos iterativos de eliminación de ruido, superan las limitaciones del VAE/GAN en cuanto a calidad de generación, estabilidad y diversidad. Sus ideas centrales incluyen: construir los datos de entrenamiento mediante la adición de ruido hacia adelante, aprender el proceso de generación mediante la eliminación de ruido hacia atrás y, finalmente, muestrear muestras de alta calidad a partir de ruido puro. Este paradigma se ha aplicado ampliamente en campos como Text-to-Image, la generación de video y el diseño molecular.

%% =====================================================
